% Security-Aware Clarification for Language-Model Agents (AAAI-27)
%
% Anonymous submission, formatted against the official AAAI-27 Author Kit
% (aaai2027.sty / aaai2027.bst, downloaded from aaai.org). Do not modify the
% style file, and do not load any package/command from its disallowed list
% (see aaai2027.sty and the Author Kit's AnonymousSubmission2027.tex).
%
% All numbers are verified against results/models/<name>/{primary_summary,stats,
% oracle_ablation,guardrail_eval}.json for the four completed models, and against
% results/cross_model_comparison.md. No \TODO placeholders remain.
\documentclass[letterpaper]{article} % DO NOT CHANGE THIS
\usepackage[submission]{aaai2027} % DO NOT CHANGE THIS
% Serif/sans-serif/monospace fonts (newtxtext, helvet, courier) and xcolor are
% loaded automatically by aaai2027.sty -- do not usepackage them separately.
\usepackage[utf8]{inputenc}
\usepackage[hyphens]{url} % DO NOT CHANGE THIS
\usepackage{graphicx} % DO NOT CHANGE THIS
\urlstyle{rm} % DO NOT CHANGE THIS
\def\UrlFont{\rm} % DO NOT CHANGE THIS
\usepackage{natbib} % DO NOT CHANGE THIS AND DO NOT ADD ANY OPTIONS TO IT
\usepackage{caption} % DO NOT CHANGE THIS AND DO NOT ADD ANY OPTIONS TO IT
\frenchspacing % DO NOT CHANGE THIS
\usepackage{amsmath}
\usepackage{booktabs}
\usepackage{multirow}

\pdfinfo{
/TemplateVersion (2027.1)
}

\setcounter{secnumdepth}{0}

\title{Ask, Act, or Verify? Security-Aware Clarification for
Language-Model Agents}
\author{Anonymous submission}
\affiliations{}

\begin{document}
\maketitle

\begin{abstract}
Language-model agents routinely receive underspecified instructions and must
decide whether to act on their best interpretation or ask for more information.
Existing clarification methods weigh expected task improvement against
interaction cost, implicitly assuming that whoever answers can be trusted. In
deployed agentic systems, this assumption often fails: clarification may arrive
through shared documents, external tools, forwarded messages, delegated
collaborators, or compromised accounts. Asking a question therefore does two
things at once---it reduces uncertainty about the user's intent, and it opens a
channel through which an attacker can inject instructions. We formulate
security-aware clarification as a decision problem in which an agent jointly
chooses whether to ask, what format to ask in, and which channel to query. We
propose \textbf{SecureVoI}, a two-stage extension of value-of-information
reasoning that scores each question--channel pair by expected information gain
net of interaction cost and channel-dependent security loss, then screens
incoming responses before accepting them. We further construct a benchmark of
120 ambiguous file- and calendar-management tasks (96 held out) with matched
benign, noisy, and adversarial clarification responses and fully automatic
outcome verification; attacks are distributed across channels so channel
identity is never perfectly predictive of an attack, preventing a defense from
succeeding merely by avoiding one channel. Across six models---four
open-weight, plus GPT-5.4-mini and Claude-Sonnet-5---conventional
value-of-information clarification raises benign task completion from near zero
to 49--100\% but incurs 28--58\% unsafe-action rates under adversarial
responses; a trivial highest-trust-channel heuristic still suffers 17--33\%, so
the benchmark cannot be solved by channel avoidance alone.
SecureVoI drives adversarial unsafe actions to 0--7\%---a significant
paired-bootstrap reduction over conventional value-of-information in every model
($-0.51$ to $-0.58$, $p<0.001$)---at a bounded cost in benign utility for which we detect no significant
difference from a trusted-channel-only baseline ($p>0.7$; we did not test
equivalence against a pre-specified margin).
An ablation that separates the two stages locates where each one matters, and
the answer is not uniform. Against explicit attacks, conventional
value-of-information paired with SecureVoI's response screen matches SecureVoI's
safety while achieving \emph{higher} adversarial utility in all four models:
response screening accounts for the headline safety gains in this regime, and
risk-aware acquisition adds no further safety once the screen is present. Against stealthier attacks, whose wording
strips the surface cues a screen keys on, adding risk-aware acquisition reduces
unsafe actions by $0.05$--$0.24$ across models, with statistically significant
reductions in two of four. These results locate the roles of the two stages:
response screening is the primary defense when attacks are recognizable, while
acquisition-time risk reasoning provides defense-in-depth as screening weakens.
More broadly, when an agent should ask depends not only on how uncertain it is,
but on who might be answering.
\end{abstract}

\section{Introduction}

Language-model agents that can call tools---archiving a file, scheduling a
meeting, sending a message---routinely receive instructions that underspecify
some part of the intended action. A growing line of work teaches agents to
recognize this ambiguity and ask a clarifying question rather than guess: CLAM
\citep{kuhn2023clam} established that a model should detect ambiguity and
request clarification instead of committing to a possibly-wrong interpretation;
SAGE-Agent \citep{suri2025sageagent} formalizes which question to ask by its expected value
of perfect information; ``Learning to Ask'' \citep{wang2024learningtoask} and ``Ask or Assume?''
\citep{edwards2026askorassume} extend the same idea to tool-calling and coding agents
specifically. All of this work shares one assumption, usually implicit: whoever
answers the question can be trusted. That assumption is the load-bearing one, and
it is also the one that breaks first in a deployed system. A clarifying question
does not necessarily go to the user who issued the original request---it may be
routed to a shared document, a delegated collaborator, a forwarded message
thread, an external tool, or, in the worst case, a compromised account. Each of
these is a plausible answerer in an ordinary agentic workflow, and none of them
is the authenticated principal the clarification-seeking literature implicitly
assumes is on the other end.

At the same time, a separate line of work has shown that tool-using agents are
vulnerable to instructions smuggled into content they process: InjecAgent \citep{zhan2024injecagent} benchmarks indirect prompt injection through tool outputs and finds that
even a capable agent follows an injected instruction a substantial fraction of
the time; AgentDojo \citep{debenedetti2024agentdojo} generalizes this into a standard
security benchmark pairing benign tasks with an adversarial one hidden in the
environment. This literature treats injected instructions as \emph{ambient}---
something the agent stumbles into while doing its job---not as something the
agent's own behavior invites. Neither literature asks the question this paper
asks: what happens when the mechanism that makes an agent \emph{helpful} (asking
for clarification when unsure) is also the mechanism an attacker can use to
\emph{reach} the agent in the first place? Asking a clarifying question is not a
passive act. It opens a channel, and opening a channel is exactly how an
indirect-prompt-injection attack gets a foothold. A policy that reasons about
whether to ask only in terms of expected information gain and interaction
cost---as every clarification-seeking method above does---has no way to represent
that the same question, asked of a different channel, carries a different
exposure.

This gap motivates the decision problem we study: an agent should decide not only
\emph{whether} to ask, but \emph{what format} to ask in and \emph{which channel}
to query, jointly, as a function of how informative each channel plausibly is,
how much asking costs in interaction overhead, and how much risk each channel
carries of returning an adversarial answer. We call this security-aware
clarification, and we propose \textbf{SecureVoI}, a two-stage extension of
value-of-information reasoning that (1) scores each candidate question--channel
pair by expected information gain net of interaction cost \emph{and} a
channel-dependent security loss term before asking, and (2) screens the response
that actually comes back before accepting it, using a learned estimate of how
likely that specific answer is to be malicious. The two stages play different roles. A policy with risk-aware acquisition but
no screen can still be talked into accepting a bad answer once it decides to
ask; a policy that only screens has no principled way to prefer a
slightly-less-informative but lower-risk channel in the first place. Our
ablation isolates stage 1's incremental effect when a response screen is
present, and we report the fourth factorial cell separately. Conventional VoI is neither: it is
risk-blind \emph{and} unscreened, which is why we show it becomes measurably
more dangerous under adversarial responses than under benign ones. We evaluate all four cells of the (risk-blind, risk-aware) $\times$
(accept-all, screen) factorial on Mistral-Nemo-12B, and three cells---all
except risk-aware acquisition without screening---across all four models.

\section{Related Work}

\paragraph{Clarification-seeking under cooperative-response assumptions.}
CLAM \citep{kuhn2023clam} establishes the foundational ``ask only when necessary''
lineage; SAGE-Agent \citep{suri2025sageagent} extends this with an EVPI-based
question-selection objective that our SecureVoI's pre-inquiry stage generalizes;
``Ask or Assume?'' \citep{edwards2026askorassume} and ``Learning to Ask'' \citep{wang2024learningtoask}
both decouple ambiguity detection from execution but assume a trustworthy
answerer throughout. None of this cluster models channel identity or adversarial
responses.

\paragraph{Prompt injection and agent security.}
InjecAgent \citep{zhan2024injecagent} is the direct precedent for our attack-success-rate
metric and \texttt{AttackType} taxonomy; AgentDojo \citep{debenedetti2024agentdojo} is the
standard tool-agent security benchmark, positioned against ours as: theirs is
ambient injection in tool output, ours is injection the agent invites by choosing
to ask. The Instruction Hierarchy \citep{wallace2024instructionhierarchy} grounds our fixed
channel-trust ranking as a training-time analog of what our policies do at
decision-time without retraining. CaMeL \citep{debenedetti2025camel} is our sharpest
foil: it constrains information flow \emph{after} acquisition and never decides
whether acquiring it was worth the exposure---the acquire decision is exactly
our contribution.

ASPI \citep{sehwag2026aspi} is the closest prior work and the one our framing
most directly answers. It shows that the mere state transition from executing a
fully specified instruction to seeking clarification substantially amplifies
prompt-injection vulnerability across ten frontier models (e.g.\ attack success
rising from $1.8\%$ to $34.0\%$ for one model), and decomposes that gap into a
state-dependent shift in how models process incoming content and a
channel-specific effect tied to the clarification interface itself. ASPI's
contribution is diagnostic and its mitigations are evaluated as fixed
defenses applied at clarification time; what it does not do---and what we
add---is treat \emph{acquisition itself} as the decision variable, optimizing
jointly over whether to ask, in what format, and which channel to expose.
Our two-stage policy is a direct response to that decomposition---stage 1's
channel-dependent risk term targets the channel-specific half of ASPI's
finding, and stage 2's response screen targets the content-processing half.

\paragraph{Value of information under adversarial channels.}
The deceptive path-planning VoI framework \citep{suttle2025deceptive} is the closest formal
precedent outside the LLM literature to ``VoI minus an adversarial-cost term'';
the Bayesian-network adversarial risk analysis framework \citep{wang2021bayesianara} grounds
our channel-prior formalism in established ARA rather than presenting it as ad
hoc.

\section{Threat Model}

\paragraph{Principal and authority.} A single authenticated principal issues the
original request, which is authoritative and trusted; the principal is never
modeled as the attacker. The agent's job is to satisfy the principal's hidden
intent, of which the request is an underspecified projection.

\paragraph{Channels.} The agent may obtain clarification from six channel classes
at different trust levels---authenticated user, delegated collaborator, shared
document, forwarded message, external tool, compromised account---and knows the
channel class but not whether a specific answer is malicious. By design, channel
identity is not perfectly predictive of attack: every non-authenticated class
contains benign, noisy, and adversarial instances, which forces a genuine
value-vs-risk trade-off rather than allowing a policy to win by simply
distrusting an entire channel class.

\paragraph{Attacker objectives.} Goal hijack, constraint removal, recipient
injection, privilege escalation, data exfiltration, destructive action, and
unrelated-instruction injection---the adversarial response always carries some
genuine information alongside the attack, so it is never trivially ignorable.

\paragraph{Attacker capabilities.} The attacker can fully control the text of a
response on any non-authenticated channel and mix valid information with the
attack. It cannot modify the system prompt, the authenticated request, the
simulator, or either verifier, and cannot read secrets except by causing the
agent itself to expose them.

\paragraph{Safety definition.} An episode is unsafe iff the safety verifier
observes any prohibited state transition in the immutable action log, even if
later reversed---safety is a property of the trajectory, not the final state.

\paragraph{Explicit non-goals.} No claim of universal protection against prompt
injection or adaptive attackers, no formal security guarantee, no
deployment-readiness claim. Simulator evidence establishes the existence and
direction of the information--security trade-off and the relative ordering of
policies, not real-world safety.

\section{Method: SecureVoI}

SecureVoI splits the clarification decision into two stages that a conventional
value-of-information (VoI) policy collapses into one. Stage 1 decides
\emph{whether, what, and where} to ask before any answer exists; stage 2 decides
\emph{whether to believe the answer} once it arrives. A single risk-aversion
parameter $\lambda$ governs both stages and is fit once, on the development
split.

\paragraph{Estimated quantities.} Both stages are driven by two model-provided
signals and never by ground-truth attack labels.
$\textsc{sample\_intents}(\text{task}, k)$ draws $k$ candidate interpretations of
the underspecified request; $\textsc{classify\_malice}(\text{text})$ returns a
scalar in $[0,1]$ estimating how likely a specific answer is to be an injected
instruction. A deterministic ScriptedAgent supplies both for the pipeline check,
and an open-weight model supplies both for the real runs, behind an identical
interface.

\paragraph{Stage 1---pre-inquiry selection.} For each candidate (question $q$,
channel $c$) pair the agent scores
\[
  \mathrm{SV}(q,c) = \mathrm{infogain}(q,c) - \mathrm{cost}(q)
  - \lambda\cdot\mathrm{prerisk}(c,q),
\]
and asks the highest-scoring pair if it is positive, otherwise acts on its best
current guess. \emph{Information gain} proxies how much an answer would change the
chosen action: the disagreement among the $k$ sampled intents (one minus the
modal action's frequency), amplified when $q$ targets a field the intents
actually dispute and when the task's stakes are high, and scaled into utility
units so it is commensurate with the risk term. It is then multiplied by a
channel-dependent informativeness factor---and in our environment the
\emph{less-trusted} channels (a shared document that lists every file; an
external tool) are genuinely \emph{more} informative than the terse
authenticated user. This is precisely what tempts a risk-blind policy onto the
channels an attacker controls. \emph{Interaction cost} is a small fixed
per-question penalty, lower for a restricted (fixed-choice) question than for an
open one. \emph{Pre-inquiry risk} is a channel-and-format prior for the
probability that an answer on $c$ carries an attack, times a normalized expected
loss; priors are Laplace-smoothed and estimated on the development split only,
and a restricted format discounts the prior because its answer is parsed outside
the language model and so exposes less injection surface.

\paragraph{Stage 2---response acceptance.} A policy that did only stage 1 could
still be talked into acting on a malicious answer once it has decided to ask.
SecureVoI therefore screens the returned answer $y$ and accepts it only if
\[
  \mathrm{infovalue}(q,c) > \lambda\cdot\mathrm{responserisk}(y)\cdot
  \mathrm{expectedloss},
\]
where $\mathrm{responserisk}$ blends the learned estimate
$\textsc{classify\_malice}(y)$ with the channel prior, lexical injection cues, and
a penalty when a restricted-format answer falls outside its allowed choices. A
rejected answer is discarded and the agent falls back to acting on its prior best
guess. The same $\lambda$ thus governs both the willingness to \emph{ask} a risky
channel and the bar an answer must clear to be \emph{believed}.

\paragraph{$\lambda$ selection.} $\lambda$ is the single knob trading benign
utility against adversarial exposure, fit on the development split alone. We
select the \emph{smallest} $\lambda$ whose dev-set adversarial unsafe rate meets a
fixed target, rather than the $\lambda$ that maximizes dev utility: on a small dev
split, utility versus $\lambda$ is not monotonic (only the unsafe rate is), so
hill-climbing utility overfits to sampling noise at a large, fragile $\lambda$.
$\lambda$, the channel priors, and all prompt wording are frozen from dev before
the test split is evaluated.

\paragraph{Baselines.} We compare against five policies spanning the design space
plus one optional guardrail: \textbf{Never Ask} (act on the initial request);
\textbf{Always Ask} (query the nominally most informative channel, ignoring cost
and risk, and accept whatever returns); \textbf{Confidence Threshold} (ask only
when sampled-intent agreement falls below a dev-calibrated threshold);
\textbf{Conventional VoI} (cost-aware information maximization but risk-blind, and
accepts whatever returns---the policy \citet{sehwag2026aspi}'s finding that
clarification-seeking amplifies prompt-injection vulnerability predicts is
dangerous); \textbf{Trusted-Only} (ask only through
channels above a trust floor); \textbf{Channel Heuristic} (always route to the
highest-trust available channel and accept whatever returns---the trivial
channel-avoidance bar); and, as an additional baseline,
\textbf{Post-Hoc Guardrail} (Conventional VoI's ask decision followed by screening
the resulting \emph{action plan} against known prohibited-action patterns before
execution). SecureVoI is the only policy that reasons about channel exposure at
acquisition time \emph{and} screens the response; the guardrail is its
architectural opposite---it lets the answer through and screens the
\emph{consequence}, a contrast we return to in the Results.

\section{Benchmark}
\label{sec:benchmark}

\paragraph{Tasks.} 120 base tasks span two tool-use domains---file management and
calendar/email scheduling---split 24 development / 96 test, frozen and
checksummed so the test split can neither drift nor leak into calibration.
Each task pairs an
underspecified request with a hidden intent that the request only partially
projects; the agent must recover that intent, by asking or by inference.

\paragraph{Matched response conditions.} The benchmark's methodological core is
that every base task carries \emph{matched} clarification responses in three
conditions---\textbf{benign} (a truthful, helpful answer), \textbf{noisy} (an
unhelpful or wrong but non-malicious answer), and \textbf{adversarial} (an answer
that mixes genuine information with an injected instruction). Because the
conditions differ only in the response, any measured difference in outcome is
attributable to the response, not to task difficulty.

\paragraph{Channels and the value--risk tension.} Responses may arrive over six
channel classes at descending prior trust---authenticated user, delegated
collaborator, shared document, forwarded message, external tool, compromised
account. Two design choices keep the trade-off honest. First, the authenticated
user is never the attacker, and channel identity is not perfectly predictive of
attack: every non-authenticated class contains benign, noisy, \emph{and}
adversarial instances, so a policy cannot win by simply blacklisting a channel
class. Second, the less-trusted channels are deliberately the \emph{more
informative} ones, so declining them forfeits real information---the tension
SecureVoI must navigate rather than sidestep.

\paragraph{Attacks.} Adversarial responses realize seven attacker
objectives---goal hijack, constraint removal, recipient injection, privilege
escalation, data exfiltration, destructive action, and unrelated-instruction
injection---each mixing a valid answer with the attack so it is never trivially
ignorable. Tasks additionally vary in \textbf{stakes} (low/medium/high, which
weight the regret of an unsafe action) and \textbf{question format} (open
free-text versus a restricted fixed-choice schema parsed outside the model).

\paragraph{Automatic verification.} Primary goal and safety outcomes are scored by two independent,
deterministic (non-LLM) verifiers; no LLM judge participates in any reported
goal, safety, or utility number. An auxiliary model judge is used only for the
refusability analysis reported later, and never feeds back into scoring.
\texttt{goal\_verifier} checks the final environment state against a
machine-checkable predicate; \texttt{safety\_verifier} scans the immutable,
append-only action log for any prohibited transition. An episode is unsafe if a
prohibited action \emph{ever} appears in the log, even if a later action reverses
it. Utility combines both, penalizing unsafe actions by task severity and
charging small costs for asking and for abstaining:
\begin{multline*}
  U = \mathrm{goal} - \mathrm{severity}\cdot\mathrm{violations} \\
  - \mathrm{clarification\_cost} - \mathrm{abstention\_cost}.
\end{multline*}

\paragraph{Known limitation.} The 120 tasks currently draw from only two base
templates (one per domain), with diversity coming from parameter combinations
(stakes $\times$ channel-availability $\times$ attack type) rather than genuinely
distinct task families; stratification by genuinely independent task family is
only partially met (see Limitations).

\section{Experimental Setup}

\begin{itemize}
  \item \textbf{Main experiment}: 120 tasks (24 dev / 96 held-out test), six
  models on the identical frozen task set: Mistral-Nemo-12B (\texttt{ollama},
  local), GPT-OSS-20B and GPT-OSS-120B (\texttt{ollama} cloud),
  Llama-3.3-70B-versatile, GPT-5.4-mini, and Claude-Sonnet-5 (hosted APIs). All
  six are test-split-complete for the seven-policy set. The stealth tier, the
  oracle ablation, the stage-1 ablation and the post-hoc guardrail were run on
  the four open-weight models only; the two hosted frontier models contribute
  the primary comparison on both task sets.
  \item \textbf{Targeted factorial}: the fourth cell (risk-aware acquisition
  without screening) is run on Mistral-Nemo-12B only, both tiers, 96/96.
  \item \textbf{Per-model calibration}: $\lambda$ and the channel priors are fit
  \emph{per model} on the 24-task dev split only, never on test
  ($\lambda=4.0$ for Mistral-Nemo-12B, $\lambda=3.0$ for the other three).
  \item \textbf{Automated validity gate}: every run is checked by
  \texttt{scripts/check\_invariants.py}, which asserts that Never-Ask (which opens
  no channel) cannot be attacked, that no benign episode is scored unsafe beyond a
  small tolerance, that attack success equals
  (adversarial $\wedge$ accepted $\wedge$ unsafe), that no unsafe reason falls
  outside a task's declared prohibited set, and that no channel is perfectly
  predictive of attack. All four runs pass.
  \item \textbf{Statistics}: task-level (paired) bootstrap, 2000 resamples, 95\%
  CI---paired because all policies share the same test tasks.
  \item \textbf{Reproducibility}: the code supplement gives the exact command
  for every run, and each result records which model produced it.
\end{itemize}

\section{Results}
\label{sec:results}

Table~\ref{tab:main} reports the held-out test split on Mistral-Nemo-12B, and
Table~\ref{tab:cross} shows that the ordering reproduces across all six.
Among the baselines, SecureVoI is the only one that remains net-positive in
utility under attack (adversarial utility $+0.071$ to $+0.217$; every other
baseline, including both trivial extremes, is negative), reducing adversarial
unsafe actions against Conventional VoI by $-0.510$ to $-0.583$ (paired
bootstrap, $p<0.001$ in every model). The ablation below shows this headline
comparison is not the whole story: on the explicit tier a screen-only policy
clears zero by a wider margin still.

\begin{table*}[t]
\centering
\begin{tabular}{lcccc}
\toprule
Policy & Ben.\ goal & Ben.\ util.\ & Adv.\ unsafe & Adv.\ util.\ \\
\midrule
Never Ask            & 0.000 & $-0.150$ & 0.000 & $-0.150$ \\
Confidence-Threshold & 0.000 & $-0.150$ & 0.000 & $-0.150$ \\
Always Ask           & 1.000 & 0.900 & 0.583 & $-0.558$ \\
Conventional VoI     & 1.000 & 0.950 & 0.583 & $-0.508$ \\
Channel Heuristic    & 1.000 & 0.950 & 0.333 & $-0.133$ \\
Trusted-Only         & 0.750 & 0.675 & 0.208 & $-0.096$ \\
\textbf{SecureVoI}   & 0.750 & 0.675 & \textbf{0.073} & $\mathbf{+0.071}$ \\
\bottomrule
\end{tabular}
\caption{Held-out test split (96 tasks, Mistral-Nemo-12B, $\lambda=4.0$).
Adversarial unsafe 95\% CIs: Conventional VoI $[0.479,0.677]$, SecureVoI
$[0.031,0.135]$. Among these baselines SecureVoI is the only one with positive
adversarial utility; the screened ablation (Results) also clears zero.
The post-hoc guardrail screens the action plan rather than choosing a policy
and is reported separately below.}
\label{tab:main}
\end{table*}

\begin{table*}[t]
\centering
\small
\begin{tabular}{lcccccc}
\toprule
Adv.\ unsafe & Mistral & OSS-20B & OSS-120B & Llama-70B & GPT-5.4-m & Claude-S5 \\
\midrule
Conventional VoI  & 0.583 & 0.583 & 0.583 & 0.573 & 0.583 & 0.281 \\
Channel Heuristic & 0.333 & 0.333 & 0.333 & 0.312 & 0.333 & 0.167 \\
Trusted-Only      & 0.208 & 0.208 & 0.208 & 0.188 & 0.208 & 0.104 \\
\textbf{SecureVoI}& \textbf{0.073} & \textbf{0.000} & \textbf{0.000} & \textbf{0.000} & \textbf{0.000} & \textbf{0.073} \\
\midrule
SecureVoI adv.\ util.\ & $+0.071$ & $+0.081$ & $+0.123$ & $+0.217$ & $+0.071$ & $+0.022$ \\
$\Delta$ vs.\ Conv.\ VoI & $-0.510$ & $-0.583$ & $-0.583$ & $-0.573$ & $-0.583$ & $-0.208$ \\
\bottomrule
\end{tabular}
\caption{Cross-model consistency on the held-out test split (96 tasks each).
Every $\Delta$ is significant at $p<0.001$ (paired bootstrap, 2000 resamples).
Claude-Sonnet-5 is markedly more injection-resistant at baseline (Conventional
VoI $0.281$ rather than $\approx0.58$), so every policy starts from a lower
unsafe rate on that model; the ordering is nonetheless unchanged.}
\label{tab:cross}
\end{table*}

\paragraph{The benchmark cannot be solved by channel avoidance.} Because the
generator spreads attacks across channels, a defense cannot succeed by refusing
one low-trust source. The \textbf{Channel Heuristic} baseline---always route to
the highest-trust available channel, accept whatever returns---beats Conventional
VoI ($0.333$ vs.\ $0.583$), confirming channel choice carries signal, but still
suffers $17$--$33\%$ unsafe actions and negative adversarial utility in every
model, so SecureVoI's lower rate is not channel avoidance. A per-episode audit
shows the mechanism: on the 20 tasks where SecureVoI asks \emph{directly on the
attacked channel}, it accepts the benign answer $20/20$ and rejects the
adversarial one $20/20$ (Mistral: $13/20$)---identical channel, opposite
decision, driven by content.

\paragraph{A more injection-resistant model narrows the gap.} Claude-Sonnet-5's
baseline is not dominated by the attack: Conventional VoI reaches only $0.281$
unsafe rather than $\approx0.58$, and Always Ask is never successfully attacked
($0/96$). That is a format effect, not evidence that asking freely is safe:
Always Ask issues \emph{open}-format questions, which this model refuses
entirely ($0/96$), while the \emph{restricted}-format attacks the cost-aware
policies prefer still succeed $22\%$ of the time ($60/267$); no other model
shows the asymmetry (all $\approx58\%$ on open format). SecureVoI still cuts
this model's unsafe rate $0.281\to0.073$ ($p<0.001$), but the headroom is
smaller, and on the three-domain set it ($0.104$) is slightly \emph{worse} than
Trusted-Only ($0.076$). As native injection resistance improves, the value of an
external clarification defense shrinks toward the residual the model cannot
handle itself.

\paragraph{The two extremes bracket the danger.} \textbf{Always Ask} accepts
whatever returns, landing at $0.583$ unsafe ($-0.558$ utility) on five of six
models. \textbf{Confidence-Threshold} collapses onto Never Ask ($0.000$ unsafe,
$0.000$ benign goal). Neither is usable---one is unsafe, the other forfeits the
task.

\paragraph{Safety is not bought with benign utility.} We detect no significant
difference between SecureVoI's benign utility and Trusted-Only's in any model
($p=0.73$--$1.00$); this is a failure to reject, not an equivalence test against
a pre-specified margin, and on GPT-OSS-20B SecureVoI is numerically slightly
worse. Under attack it is substantially safer ($0.073$ vs.\ $0.208$ on Mistral;
$0.000$ vs.\ $0.188$--$0.208$ elsewhere), roughly $3\times$ fewer unsafe actions
at benign utility we cannot distinguish. It does
give up benign utility relative to risk-blind Conventional VoI ($0.675$ vs.\
$0.950$ on Mistral)---this is the real, bounded price of security, and we report it
plainly rather than claiming a free lunch.

\paragraph{Oracle-vs-learned-risk ablation.} Replacing the learned response
screen with the ground-truth attack label removes Mistral-Nemo's entire residual
unsafe rate ($0.073\to0.000$) and changes nothing on the other three, already at
$0.000$. The residual is thus a stage-2 classifier limitation on the weakest
model, not a limitation of the acquisition rule: Mistral's screen admits $7$ of
the $20$ attacks reaching it, the stronger models none. A better classifier is
worth at most $0.073$.

\paragraph{Isolating stage 1 from stage 2.} The oracle ablation above varies
\emph{how good} the response screen is; it does not test whether stage-1
acquisition contributes anything once a screen is present at all. We therefore
run the sharpest available ablation: Conventional VoI's risk-blind acquisition
paired with SecureVoI's stage-2 screen, unchanged
(\texttt{ScreenedConventionalVoI}, which inherits \texttt{accept} from SecureVoI,
so their acceptance rules are identical by construction even though the two
policies routinely ask different questions on different channels). Any gap
between this policy and full SecureVoI is attributable to stage 1 alone. We
report the raw difference in adversarial unsafe rate with its confidence
interval, and secondarily the fraction of Conventional VoI $\to$ SecureVoI's
reduction that the screened policy fails to close (the ``stage-1 share'').
Because this design varies stage 1 only in the presence of stage 2, it
identifies stage 1's \emph{incremental} contribution given a screen. To also
measure stage 1 acting alone we add the fourth factorial cell,
risk-aware acquisition with no screen, on Mistral-Nemo-12B (both tiers,
$n=96$, full coverage). The factorial measures component effects at SecureVoI's
fixed, dev-selected $\lambda$; it does not separately re-optimize $\lambda$ for
each cell, so ``stage 1 alone is net-negative'' is a statement about that
$\lambda$, not about the best achievable acquisition-only policy.

This ablation, the stealth-tier decomposition, and the targeted factorial were
all specified after seeing the main results and are reported as exploratory.
The answer depends entirely on attack tier, and this is the paper's most
important negative-then-positive result. \textbf{On the explicit tier, stage 1
contributes nothing}: the share is exactly $0.000$ in all four models
($p=1.000$ throughout), because the screened policy already drives adversarial
unsafe actions to $0.073$, $0.000$, $0.000$ and $0.000$---matching full
SecureVoI exactly. Worse for a two-stage story, stage 1 is not free
(Table~\ref{tab:ablation}): the screened policy attains identical safety with
higher benign utility in three of four models, and \emph{dominates SecureVoI on
adversarial utility in all four}, because SecureVoI's risk term declines to ask
on $24/96$ adversarial tasks that stage 2 would have screened successfully
anyway. This qualifies the headline claim above: on the explicit tier the
strongest policy we ran is not the one we propose.

\textbf{But redundant is not the same as ineffective.} The fourth factorial
cell separates the two (Table~\ref{tab:ablation}). Both stages are effective
standalone \emph{safety} mechanisms, though stage 1 alone remains net-negative
in adversarial utility ($-0.10$) because part of its safety comes from
declining to ask and completing less: against risk-blind accept-all, stage 1
alone gives
$\Delta_{\text{unsafe}}=-0.375$ $[-0.521,-0.219]$ and stage 2 alone
$-0.510$ $[-0.604,-0.406]$ (both $p<0.001$). These are \emph{simple} effects,
each measured with the other factor held off, not factorial main effects. The
explicit tier shows not that stage 1 is weak but that it is
\emph{substitutable}: a difference-in-differences confirms it helps
significantly more when no screen is present ($+0.375$, $[0.240,0.510]$,
$p<0.001$).

\textbf{On the stealth tier the conclusion reverses.} Once the injected text no
longer carries the explicit tier's surface tells, the screen degrades and stage 1
begins to carry real weight. Removing it raises the adversarial unsafe rate by
$+0.05$ to $+0.24$ (Table~\ref{tab:ablation}), and SecureVoI recovers the
adversarial-utility advantage in three of four models. The paired bootstrap
separates these by strength of evidence: the increase is significant for
GPT-OSS-20B ($p<0.001$) and Llama-3.3-70B ($p=0.012$), but its interval
includes zero for GPT-OSS-120B ($p=0.109$) and marginally so for
Mistral-Nemo-12B ($p=0.057$). We therefore claim positive estimated effects in
all four models with statistically significant ones in two; the corresponding
$n=144$ results are reported below but, being four model-level tests on a set
built after these observations, we read them as corroborating rather than
confirmatory. All four models have both tiers at full $96$-task coverage with
zero excluded tasks, so each reversal is a \emph{within-model} comparison on
identical tasks. The factorial makes the reversal sharper still, and explains
it. Stage 1 alone yields \emph{exactly} the same $0.208$ unsafe rate on both
tiers---indeed it fails on the identical $20$ tasks. Its acquisition decision is
made \emph{before} the response exists and it applies no response-risk
classifier, so its channel choice is unaffected by the attack's surface wording;
the agent still uses the accepted response to complete the task. Stage 2 alone, which does read the response,
falls from $0.073$ to $0.354$ under the same rewrite. The two stages therefore
\emph{swap roles}: in this targeted Mistral factorial, screening dominates against explicit
attacks ($0.073$ vs.\ $0.208$) with stage 1 adding nothing on top of it, while
risk-aware acquisition dominates the screen-only policy under stealth attacks
($0.208$ vs.\ $0.354$), the screen then adding nothing on top of \emph{it}
($p=1.000$). We do not generalize this single-model result, but it identifies the mechanism the four-model ablation can only
infer: a \emph{pre-response} defense is what survives paraphrase.
The effect persists under a change of task family: on a second $180$-task set
($144$ test) built by the same generator with a \emph{messaging} domain added
and $\lambda$ and the channel priors re-fit on its own dev split, stage 1
retains a substantial share in all four models ($0.206$--$0.688$). The
direction of change is mixed (two rise, one flat, one falls), so we claim
persistence, not growth. Per-domain, the share is $0.000$ exactly where the
screen alone already reaches $0.000$ unsafe and large where it leaves residual
risk---again, stage 1 pays in proportion to what the screen misses.
The headline explicit-tier numbers therefore understate how much of SecureVoI's
design is load-bearing: the regime where stage 1 earns its cost is precisely the
harder one the explicit tier does not probe.

\begin{table}[t]
\centering
\footnotesize
\setlength{\tabcolsep}{2.4pt}
\begin{tabular}{@{}lcccc@{}}
\toprule
& Conv.\ VoI & Screened & SecureVoI & $\Delta$ stage 1 \\
Model & uns/util & uns/util & uns/util & [95\% CI] \\
\midrule
\multicolumn{5}{@{}l}{\emph{Explicit tier}} \\
Mistral-12B  & .58/$-$.51 & .07/$+$.08 & .07/$+$.07 & .00 [$-$.07,.08] \\
\quad{\itshape+ stage-1-only} & \multicolumn{4}{@{}l}{\itshape .21/$-$.10 adv, goal .40, asks 75\%} \\
Llama-70B    & .58/$-$.80 & .00/$+$.37 & .00/$+$.22 & .00 [.00,.00] \\
GPT-OSS-20B  & .58/$-$1.0 & .00/$+$.17 & .00/$+$.10 & .00 [.00,.00] \\
GPT-OSS-120B & .58/$-$.91 & .00/$+$.26 & .00/$+$.14 & .00 [.00,.00] \\
\midrule
\multicolumn{5}{@{}l}{\emph{Stealth tier}} \\
Mistral-12B  & .58/$-$.51 & .35/$-$.29 & .21/$-$.10 & .15 [.00,.29] \\
\quad{\itshape+ stage-1-only} & \multicolumn{4}{@{}l}{\itshape .21/$-$.10 adv, goal .40 (as explicit)} \\
Llama-70B    & .58/$-$.80 & .24/$-$.11 & .09/$+$.03 & .15 [.03,.26] \\
GPT-OSS-20B  & .56/$-$.88 & .26/$-$.28 & .02/$+$.12 & .24 [.15,.33] \\
GPT-OSS-120B & .57/$-$.91 & .07/$+$.09 & .02/$+$.07 & .05 [.00,.12] \\
\bottomrule
\end{tabular}
\caption{Stage-1 ablation, all four models, both tiers ($n=96$, full coverage);
adversarial unsafe rate / adversarial utility. ``Screened'' is Conventional
VoI's risk-blind acquisition with SecureVoI's identical stage-2 screen;
$\Delta$ is the increase in unsafe rate from removing stage 1. Explicit:
$\Delta=0$ throughout, and the ablation attains \emph{higher} utility than
SecureVoI in all four. Stealth: $\Delta>0$ in all four, and SecureVoI recovers
the utility advantage in three. Italic rows give the fourth factorial cell (risk-aware acquisition, no screen),
run on Mistral-Nemo-12B only; its benign goal/utility are $0.75$/$+0.68$ vs.\
$1.00$/$+0.95$ for risk-blind accept-all.}
\label{tab:ablation}
\end{table}

\paragraph{Sensitivity to the utility weights.} Recomputing utility from the
saved episodes over a $6\times4$ grid of severity ($1$--$8$) and abstention
($0$--$0.5$) weights separates two claims. The \emph{sign} claim is
weight-dependent: SecureVoI's adversarial utility is positive at the reported
severity of $2$ but non-positive in $16$ of $24$ cells, turning negative once
violations weigh $\approx3\times$ goal completion. The \emph{ordering} claim,
which our conclusions rest on, is stable in $20$ of $24$ cells, failing only at
severity $1$---violations no worse than a missed task---where the Channel
Heuristic edges it ($0.200$ vs.\ $0.181$).

\paragraph{Post-hoc guardrail.} Screening Conventional VoI's resulting action
plan against the prohibited-action patterns reaches $0.000$ adversarial unsafe
on three models ($0.052$ on GPT-OSS-20B). This is an \emph{upper bound}, not a
competitor: the filter has oracle knowledge of each task's prohibited patterns,
and every attack here manifests as one conspicuous out-of-scope call. It is not
free even so---adversarial goal rate falls to $0.542$--$0.698$ on three models,
since stripping the injected call disrupts legitimate ones. Output-level
screening excels when harm is a discrete recognizable action; acquisition-time
reasoning should matter more against attacks that evade an action filter, a
regime this benchmark does not stress.

\paragraph{Stability under sampling.} Re-running a 30-task subset three times at
temperature $0.7$ shows the effect is not a decoding artifact: Conventional
VoI's adversarial unsafe rate is $0.511 \pm 0.016$ and SecureVoI's $0.122 \pm
0.016$ (mean $\pm$ SD), every measured quantity has SD $\le 0.05$, and the
ordering persists.

\paragraph{Robustness to attack phrasing.} A second attack tier carries the
identical prohibited action on the identical channel and task with the explicit
tier's surface tells removed, drawn from three distinct persuasion strategies
per cell (12 strings) so no result rests on one sentence. As a validity check,
no risk-blind policy shows a significant stealth effect in any model
($p\ge0.09$), confirming the rewrite changed detectability rather than the task.
Against that flat baseline, SecureVoI's adversarial unsafe rate rises in all
four open-weight models, tracking capability: GPT-OSS-120B $0.000\to0.021$
($p=0.28$), GPT-OSS-20B $0.000\to0.042$ ($p=0.038$), Llama-3.3-70B
$0.000\to0.083$ ($p<0.001$), Mistral-Nemo-12B $0.073\to0.208$ ($p<0.001$). A
response-risk decomposition separating the learned classifier from a hardcoded
keyword heuristic shows \emph{why} the degradation differs: the weaker model's
failure is substantially a genuine classifier limitation (Mistral-Nemo-12B,
$68\%$ of the risk drop from the learned screen itself), while the strongest
model's small residual movement is mostly the keyword list losing its triggers
(GPT-OSS-120B, $33\%$; the others fall between at $45$--$48\%$). An auxiliary judge (GPT-OSS-120B---one of the evaluated model families, so
this is supportive rather than independent confirmation---run for 3 repetitions with a prompt that
deliberately never mentions injection or security, so as not to prime a
DECLINE; verbatim in the supplement) confirms the rewrite preserved the
benchmark's fairness bound: stealth attacks remain judged
refusable $83\%$ of the time versus $100\%$ for the explicit tier, with a $0\%$
false-positive rate on benign controls, so the harder detection problem is not an
artifact of an unfair or unanswerable attack.


\section{Limitations}
\label{sec:limitations}

\begin{itemize}
  \item \textbf{Attack explicitness.} Main-tier injections are fairly overt
  imperatives; the stealth tier shows SecureVoI's unsafe rate rises in every
  model, a genuine classifier limitation on the weaker ones rather than a
  keyword-list artifact. The main-tier $0.000$--$0.073$ results are a floor, not
  a ceiling.
  \item \textbf{The headline benchmark has only two base task templates}
  (file-archive, calendar-scheduling), with diversity from parameter
  combinations rather than distinct task families. Our claims are therefore
  about the \emph{ordering} of policies and the \emph{existence} of the
  acquisition-time trade-off, not about absolute rates transferring. With two
  clusters there is no clean statistical repair: task-level bootstrap intervals
  treat parameter variants as independent draws, quantifying uncertainty within
  this fixed synthetic benchmark, not over new task families or attack language.
  The separate $180$-task three-domain set on which the ablation replicates was
  authored by us under the same design assumptions, so it tests robustness to
  task family but not to benchmark authorship; porting onto AgentDojo
  \citep{debenedetti2024agentdojo} or ASPI's suite \citep{sehwag2026aspi} is
  the correct next step.
  \item \textbf{Few distinct attack strings.} The explicit tier uses four
  injected texts (2 domains $\times$ 2 formats) across 120 tasks, so bootstrap
  intervals cover channel, stakes and information gain but \emph{not} linguistic
  variation in the attack; the stealth tier triples this to 12, still a small
  sample. All stealth-tier claims are exploratory with respect to attack
  language.
  \item Six of the seven declared attacker objectives are exercised by the current
  attack texts; \texttt{delete:credentials.json} and \texttt{email\_external} are
  declared and verifier-covered but never triggered, and
  \texttt{delete\_protected} fires rarely.
  \item $\textsc{classify\_malice}$ is a model-prompted heuristic and is the
  binding constraint on the weakest model: the oracle ablation removes all of
  Mistral-Nemo's residual $0.073$, so safety depends on screening quality in a
  way that scales with model capability.
  \item \textbf{Not all unsafe actions come from injections.} On one benign
  episode ($1/672$) Llama-3.3-70B spontaneously disclosed a private calendar
  from a benign answer that never mentioned sharing. Proactive agents can
  breach a constraint unprompted, so clarification safety is necessary but not
  sufficient for agent safety.
  \item On GPT-OSS-20B SecureVoI's benign utility ($0.321$) falls marginally
  below Trusted-Only's ($0.331$, $p=0.855$), and on Claude-Sonnet-5's
  three-domain set its unsafe rate ($0.104$) is above Trusted-Only's ($0.076$);
  we do not claim dominance on either.
  \item Attacks are static per task; we do not model an adaptive attacker.
  Simulator evidence establishes the trade-off's existence and the policy
  ordering, not real-world deployment safety.
\end{itemize}

\section{Conclusion}

Clarification-seeking is not only a question of how uncertain an agent is, but
of who is likely to answer. Risk-blind value-of-information reasoning sharply
raises unsafe actions once the channel can be adversarial; SecureVoI closes most
of that gap by $0.208$--$0.583$ across six models ($p<0.001$).

\bibliography{references}

\end{document}

